% Compile with: pdflatex homework_solutions.tex (twice for the contents).
% Append later assignments before \end{document}; a template is at the end.
\documentclass[11pt,a4paper]{article}
\usepackage[T1]{fontenc}
\usepackage[utf8]{inputenc}
\usepackage{lmodern}
\usepackage[margin=27mm]{geometry}
\usepackage{amsmath,amssymb,amsthm}
\usepackage[hidelinks]{hyperref}
\usepackage{microtype}
\setlength{\parindent}{0pt}
\setlength{\parskip}{6pt}
\allowdisplaybreaks

\newcounter{question}[section]
\newcommand{\homework}[1]{%
  \clearpage
  \section{Homework #1}%
  \setcounter{question}{0}%
  \setcounter{equation}{0}%
}
\renewcommand{\theequation}{\thesection.\arabic{equation}}
\newenvironment{question}{%
  \refstepcounter{question}%
  \subsection*{Question \arabic{question}}%
  \addcontentsline{toc}{subsection}{Question \arabic{question}}%
}{\par}
\newenvironment{answer}{\par\medskip\textbf{Answer.}\par}{\par\medskip}
\newcommand{\eps}{\varepsilon}
\DeclareMathOperator{\erf}{erf}
\DeclareMathOperator{\erfc}{erfc}

\title{Asymptotic Methods\\Homework Questions and Solutions}
\author{}
\date{}

\begin{document}
\maketitle
\tableofcontents

\homework{1}
\textit{Source: \texttt{homework1.pdf}, Exercise examples \#1,
Fall 2026--2027. All three numbered problems are included.
In Question 3, the source mixes $x$ and $z$ in the notation for the
error function; the variable is written consistently as $z$ below.}

\begin{question}
Show that
\[
  f\sim g \quad\Longleftrightarrow\quad f=g[1+o(1)].
\]
\end{question}

\begin{answer}
Fix a limiting process, for example $x\to a$ or $x\to\infty$.
Assume that $g(x)\ne0$ throughout a sufficiently small punctured
neighborhood of the limit (or for all sufficiently large $x$).
Both statements must refer to this same limiting process.

By definition, asymptotic equivalence means
\[
  f(x)\sim g(x)
  \quad\Longleftrightarrow\quad
  \lim_{x\to a}\frac{f(x)}{g(x)}=1.
\]
The notation $r(x)=o(1)$ means $r(x)\to0$.

\textbf{Forward implication.}
Suppose $f\sim g$ and define
\[
  r(x)=\frac{f(x)}{g(x)}-1.
\]
Then $r(x)\to0$, so $r=o(1)$, and rearranging gives
\[
  f(x)=g(x)\bigl[1+r(x)\bigr]=g(x)[1+o(1)].
\]

\textbf{Reverse implication.}
Suppose $f(x)=g(x)[1+r(x)]$ for some $r(x)\to0$.
Division by $g(x)$ gives
\[
  \frac{f(x)}{g(x)}=1+r(x)\longrightarrow1.
\]
Thus $f\sim g$, proving the equivalence.

An equivalent formulation is $f-g=o(g)$, meaning that
$(f-g)/g\to0$. The error is small \emph{relative to $g$};
the absolute difference need not tend to zero.
For example, $x+\sqrt{x}\sim x$ as $x\to\infty$, even though
their difference tends to infinity.
\end{answer}

\begin{question}
Obtain two-term approximations for the two large roots of
\[
  \eps^{1/2}z^3+(1+\eps)z+\eps^{1/3}=0.
\]
\end{question}

\begin{answer}
We take $\eps\to0^+$ and allow complex roots.
The two large roots are a complex-conjugate pair.

\textbf{1. Find the scale by dominant balance.}
For a large root, balance the cubic and linear terms:
\[
  \eps^{1/2}z^3\sim-z,
  \qquad z^2\sim-\eps^{-1/2}.
\]
Consequently, $z=O(\eps^{-1/4})$ and the leading approximations are
$z\sim\pm i\eps^{-1/4}$.
The constant term is small on this scale.

Set $z=\eps^{-1/4}w$ and multiply the equation by $\eps^{1/4}$:
\begin{equation}\label{eq:scaled}
  w^3+(1+\eps)w+\eps^{7/12}=0.
\end{equation}
At $\eps=0$, the limiting equation is $w(w^2+1)=0$.
The roots $w_0=\pm i$ give the two large roots of the original
equation. The root $w_0=0$ belongs to a different scale.

\textbf{2. Determine the first correction.}
Write $w=w_0+\delta$, with $w_0=\pm i$ and $\delta=o(1)$.
Using $w_0^3+w_0=0$ and $3w_0^2+1=-2$ in
\eqref{eq:scaled}, we obtain
\[
  -2\delta+3w_0\delta^2+\delta^3
  +\eps w_0+\eps\delta+\eps^{7/12}=0.
\]
Since $\eps=o(\eps^{7/12})$, the largest forcing is
$\eps^{7/12}$. Balancing it with $-2\delta$ gives
\[
  \delta\sim\frac12\eps^{7/12}.
\]
More precisely, substitution shows that
\[
  w_\pm=\pm i+\frac12\eps^{7/12}+O(\eps).
\]
The validity of this local expansion can also be seen by putting
$t=\eps^{1/12}$ in \eqref{eq:scaled}. The resulting equation,
$w^3+(1+t^{12})w+t^7=0$, is analytic in $(w,t)$, and its derivative
with respect to $w$ is $-2$ at $(\pm i,0)$. The implicit function
theorem therefore gives a regular local expansion in $t$.

\textbf{3. Return to the original variable.}
Multiplying by $\eps^{-1/4}$ yields the requested two-term answers:
\begin{equation}
  \boxed{z_\pm(\eps)=\pm i\eps^{-1/4}
       +\frac12\eps^{1/3}+O(\eps^{3/4}).}
\end{equation}
The correction is smaller than the leading term by a factor
$\eps^{7/12}\to0$.

As an additional check, retaining the next forcing $\eps w_0$
gives
\[
  z_\pm=\pm i\eps^{-1/4}+\frac12\eps^{1/3}
       \pm\frac{i}{2}\eps^{3/4}+O(\eps^{11/12}).
\]
The remaining, small root has leading behavior
$z_{\mathrm{small}}\sim-\eps^{1/3}$, obtained by balancing the
linear and constant terms. This also explains why there are
exactly two large roots and one small root.
\end{answer}

\begin{question}
Obtain an asymptotic series for
\[
  \erf z=\frac{2}{\sqrt{\pi}}\int_0^z e^{-t^2}\,dt
  \qquad\text{as }z\to\infty.
\]
Hint: recall that $\erf(\infty)=1$.
Does the series converge? How many terms are needed to obtain
an accuracy of $10^{-5}$ for $z=3$?
\end{question}

\begin{answer}
We consider real $z\to+\infty$ and interpret the requested
accuracy as an absolute error in $\erf(3)$.

\textbf{1. Expand the tail integral.}
Using the hint, define the complementary error function:
\[
  \erfc z=1-\erf z
  =\frac{2}{\sqrt{\pi}}\int_z^\infty e^{-t^2}\,dt.
\]
For $n\ge0$, put
\[
  I_n(z)=\int_z^\infty t^{-2n}e^{-t^2}\,dt.
\]
Since $d(e^{-t^2})=-2t e^{-t^2}\,dt$, integration by parts gives
\begin{align*}
  I_n(z)
  &=-\frac12\int_z^\infty t^{-(2n+1)}\,d(e^{-t^2})\\
  &=\frac{e^{-z^2}}{2z^{2n+1}}
     -\frac{2n+1}{2}I_{n+1}(z).
\end{align*}
The boundary term at infinity vanishes because exponential decay
dominates every fixed power.

Repeated use of this identity gives
\begin{equation}\label{eq:erfc-expansion}
  \erfc z=
  \frac{e^{-z^2}}{\sqrt{\pi}\,z}
  \sum_{n=0}^{N-1}(-1)^n\frac{(2n-1)!!}{(2z^2)^n}
  +R_N(z),
\end{equation}
where $(-1)!!=1$ and
\[
  R_N(z)=(-1)^N\frac{2}{\sqrt{\pi}}
         \frac{(2N-1)!!}{2^N}I_N(z).
\]
In particular, the asymptotic expansion is
\begin{equation}
  \boxed{\erf z\sim1-
  \frac{e^{-z^2}}{\sqrt{\pi}\,z}
  \left(1-\frac{1}{2z^2}+\frac{3}{4z^4}
        -\frac{15}{8z^6}+\frac{105}{16z^8}-\cdots\right).}
\end{equation}

\textbf{2. Bound the remainder and justify the expansion.}
For $z>0$, the recurrence and $I_{N+1}(z)>0$ imply
\[
  0<I_N(z)<\frac{e^{-z^2}}{2z^{2N+1}}.
\]
Thus
\begin{equation}\label{eq:remainder-bound}
  |R_N(z)|<\frac{e^{-z^2}}{\sqrt{\pi}\,z}
            \frac{(2N-1)!!}{(2z^2)^N}.
\end{equation}
The error is smaller than the magnitude of the first omitted
term and has that term's sign in the expansion for $\erfc z$.
For each fixed $N$, this bound is of smaller order than the last
retained term as $z\to+\infty$. This proves the asymptotic property.

\textbf{3. Does the infinite series converge?}
Let $a_n=(2n-1)!!/(2z^2)^n$ be the magnitude of the $n$th term
inside the parentheses. For any fixed finite $z>0$,
\[
  \frac{a_{n+1}}{a_n}=\frac{2n+1}{2z^2}\longrightarrow\infty.
\]
The terms eventually grow and fail to tend to zero. Therefore the
series diverges for every fixed finite $z>0$.
There is no contradiction: an asymptotic expansion describes the
error after a fixed number of terms as $z\to\infty$, whereas
convergence would require a limit as $N\to\infty$ with $z$ fixed.
In practice, the series should be truncated near its smallest term
if especially high accuracy is wanted.

\textbf{4. Number of terms for $z=3$.}
Retaining just the first term of the tail expansion gives
\[
  \erf(3)\approx1-\frac{e^{-9}}{3\sqrt{\pi}}
              \approx0.999976791158.
\]
For $N=1$, \eqref{eq:remainder-bound} gives the rigorous bound
\[
  |\text{error}|<\frac{e^{-9}}{3\sqrt{\pi}}\frac1{18}
               \approx1.28938\times10^{-6}<10^{-5}.
\]
Hence \emph{one term of the tail series} is sufficient, in addition
to the leading constant $1$ in $\erf z$.

To check that the constant $1$ alone is insufficient, use
$I_0=e^{-z^2}/(2z)-I_1/2$ and
$I_1<I_0/z^2$ to obtain
\[
  \erfc z>
  \frac{e^{-z^2}}{\sqrt{\pi}\,z}
  \frac{1}{1+1/(2z^2)}.
\]
At $z=3$, this lower bound is approximately
$2.19873\times10^{-5}>10^{-5}$.
Thus the minimal answer is
\[
  \boxed{\begin{gathered}
    \text{One tail term (two contributions including the constant):}\\
    \erf(3)\approx1-\frac{e^{-9}}{3\sqrt{\pi}}.
  \end{gathered}}
\]
\end{answer}

% =====================================================================
% FUTURE ASSIGNMENTS
% Uncomment and repeat this block for each new homework assignment.
% Each \homework command starts a new page and resets question numbering.
% Equations are numbered by assignment section.
%
% \homework{2}
% \begin{question}
% Paste the next question here.
% \end{question}
% \begin{answer}
% Write the detailed solution here.
% \end{answer}
%
% \begin{question}
% Paste another question here.
% \end{question}
% \begin{answer}
% Write its detailed solution here.
% \end{answer}
% =====================================================================

\end{document}
